% !TEX program = xelatex
% 공통수학2 · Ⅱ. 집합과 명제 · 1. 집합 · 01 집합 · 2/2차시
\documentclass[11pt,a4paper]{article}
\usepackage{kotex}
\usepackage{iftex}
\ifPDFTeX\else
  \setmainhangulfont{Noto Serif CJK KR}
  \setsanshangulfont{Noto Sans CJK KR}
\fi
\usepackage[margin=2.1cm]{geometry}
\usepackage{parskip}
\usepackage{enumitem}
\usepackage{array}
\usepackage{tabularx}
\usepackage{booktabs}
\usepackage{xcolor}
\usepackage{amssymb}
\usepackage{tikz}
\usepackage[most]{tcolorbox}
\usepackage{fancyhdr}
\usepackage{titlesec}

\definecolor{goldc}{HTML}{B07C22}
\definecolor{goldstrong}{HTML}{8A5F16}
\definecolor{indigoc}{HTML}{2B3B57}
\definecolor{indigosoft}{HTML}{6E82A6}
\definecolor{linec}{HTML}{D6DACB}
\definecolor{surface2}{HTML}{E4E8DC}
\definecolor{notebg}{HTML}{FBF3E3}
\definecolor{inksoft}{HTML}{52604F}

\titleformat{\section}{\normalfont\sffamily\bfseries\large\color{indigoc}}{\thesection}{0.6em}{}
\titlespacing{\section}{0pt}{16pt}{8pt}
\newtcolorbox{ideabox}{enhanced, breakable, colback=white, colframe=goldc, boxrule=0.5pt, leftrule=3pt, arc=2pt, left=10pt, right=10pt, top=8pt, bottom=8pt}
\newtcolorbox{calloutbox}{enhanced, breakable, colback=white, colframe=indigosoft, boxrule=0.5pt, leftrule=3pt, arc=2pt, left=10pt, right=10pt, top=7pt, bottom=7pt}
\newtcolorbox{notebox}{enhanced, breakable, colback=notebg, colframe=goldc, boxrule=0.5pt, leftrule=3pt, arc=2pt, left=10pt, right=10pt, top=7pt, bottom=7pt}
\newtcolorbox{answerbox}{enhanced, breakable, colback=white, colframe=linec, boxrule=0.5pt, arc=2pt, left=10pt, right=10pt, top=7pt, bottom=7pt}
\newcommand{\lessonbanner}[3]{%
  \begin{tcolorbox}[enhanced, colback=indigoc, colframe=indigoc, boxrule=0pt, arc=0pt,
    left=14pt, right=14pt, top=14pt, bottom=10pt, borderline south={2.4pt}{0pt}{goldc}, fontupper=\color{white}]
    {\sffamily\footnotesize\color{white!85!black} #1}\\[4pt]{\sffamily\bfseries\Large #2}\\[3pt]{\sffamily\small\color{white!88!black} #3}
  \end{tcolorbox}}
\pagestyle{fancy}\fancyhf{}\renewcommand{\headrulewidth}{0pt}
\fancyfoot[L]{\footnotesize\color{inksoft} 공통수학2 · 2026학년도 2학기 · 지평선고등학교}
\fancyfoot[R]{\footnotesize\color{inksoft} 교사용 지도안 -- 배포 금지}

\begin{document}
\lessonbanner{공통수학2 · Ⅱ. 집합과 명제 · 1. 집합}{01 집합 --- 2/2차시}{원소의 개수와 벤 다이어그램 · 교사용 지도안 (2026학년도 2학기)}
\vspace{10pt}

\noindent\begin{tabularx}{\linewidth}{@{}>{\bfseries\color{indigoc}}p{2.6cm}X@{}}
\toprule
대상 & 고1 · 2개 반 · 반당 13명 \\
차시 & 50분 (도입 10 · 전개 30 · 정리 10) \\
성취기준 & 집합의 개념을 이해하고, 집합을 표현할 수 있다. \\
선행 학습 & 18차시 --- 집합의 뜻과 표현 \\
\bottomrule
\end{tabularx}

\section{학습 목표 \& Big Idea}
\begin{ideabox}
\textbf{\color{goldstrong}영속적 이해 (Big Idea)}\\
집합은 `크기'(원소의 개수)로도 분류할 수 있다. 그런데 `아무것도 없음'(공집합)도 하나의 온전한 집합이라는 사실은, `없음' 역시 다룰 수 있는 대상이 될 수 있음을 보여 준다.
\vspace{4pt}
\small\color{inksoft} 핵심개념: \textbf{크기(개수)} \quad\textbullet\quad 핵심개념: \textbf{없음도 대상이 된다} \quad\textbullet\quad 일반화: \textbf{경계의 개념(공집합)도 엄밀하게 정의될 수 있다}
\end{ideabox}
\begin{calloutbox}
\textbf{차시 학습 목표}
\begin{itemize}[leftmargin=1.4em, itemsep=2pt, topsep=4pt]
  \item 집합을 벤 다이어그램으로 표현할 수 있다.
  \item 유한집합, 무한집합, 공집합을 구분하고, 유한집합의 원소의 개수 $n(A)$를 구할 수 있다.
\end{itemize}
\end{calloutbox}

\section{질문의 연결}
\begin{tabularx}{\linewidth}{@{}>{\bfseries\color{indigoc}\small}p{2.4cm}X@{}}
사실적 질문 & 원소가 하나도 없는 모임도 집합이라고 부를 수 있을까? \\[6pt]
개념적 질문 & 집합 $\{0\}$과 공집합 $\varnothing$은 왜 서로 다를까? \\[6pt]
논쟁적 질문 & `텅 빈 것'을 하나의 대상으로 인정하고 이름 붙이는 것 --- 수학 밖의 삶에서도 `없음'을 인정하는 것이 왜 중요할까? \\
\end{tabularx}

\section{수업 흐름}
\begin{tabularx}{\linewidth}{@{}>{\centering\arraybackslash}p{1.6cm}X>{\raggedright\arraybackslash\small}p{3.1cm}@{}}
\toprule
\textbf{단계} & \textbf{활동 \& 발문} & \textbf{ATL / VTR} \\
\midrule
\textcolor{goldstrong}{\textbf{도입}}\newline\footnotesize 10분 &
\textbf{마음공부 (2분)} --- 유무념 대조.\newline
\textbf{생각 열기 (8분)} --- 집합 $A=\{a,b,c,d,e\}$를 벤 다이어그램으로 그려 보고, 원소가 하나도 없는 집합은 어떻게 그릴지 예측해 보게 한다. &
ATL: 사고(시각화)\newline VTR: See-Think-Wonder \\
\addlinespace
\textcolor{goldstrong}{\textbf{전개}}\newline\footnotesize 30분 &
\textbf{벤 다이어그램 (10분)} --- 문제 4 짝 활동(원소 나열 집합, 조건제시 집합 각각 그리기).\newline
\textbf{원소의 개수 (20분)} --- 유한집합·무한집합 구분 $\to$ 공집합 $\varnothing$과 $n(\varnothing)=0$ 소개 $\to$ ``집합 $\{0\}$은 공집합일까?'' 발문(오개념 확인) $\to$ 문제 5(방정식의 해가 없는 경우 포함). &
ATT: 촉진자 역할\newline ATL: 개념 구별 \\
\addlinespace
\textcolor{goldstrong}{\textbf{정리}}\newline\footnotesize 10분 &
\textbf{개념 연결 질문 (5분)} --- ``$\{0\}$과 $\varnothing$이 다른 이유를 한 문장으로?''\newline
\textbf{성찰 (5분)} --- 감사 일기 1문장 + `집합' 소단원 마무리 + 다음 차시 예고(집합 사이의 포함 관계). &
마음공부 · 감사 일기 \\
\bottomrule
\end{tabularx}

\begin{calloutbox}
\textbf{지도상의 유의점} --- 집합 $\{0\}$은 원소 $0$을 가지므로 공집합이 아니다(원소가 1개인 집합)! 이는 학생들이 매우 흔하게 틀리는 오개념이므로 명시적으로 짚어 준다. 원소의 개수는 유한집합에서만 정의되며 무한집합에서는 다루지 않는다.
\end{calloutbox}

\section{형성평가 \& 답안}
\begin{answerbox}
\textbf{문제 5} \quad $n(A)$ 구하기\\
\small ⑴ $A=\{3,6,9\}$ \quad ⑵ $A=\{1^2,2^2,3^2,\dots,10^2\}$ \quad ⑶ $A=\{x\mid x^2=-4$인 실수$\}$ \quad ⑷ $A=\{x\mid x^2-2x-3=0\}$\\[4pt]
\textbf{\color{goldstrong}답} \; ⑴ $3$ \quad ⑵ $10$ \quad ⑶ $0$ (공집합) \quad ⑷ $2$
\end{answerbox}

\section{성취수준 (이 차시 범위)}
\begin{tabularx}{\linewidth}{@{}>{\bfseries\color{goldstrong}}p{0.9cm}X@{}}
\toprule
A & 집합의 개념을 설명하고, 집합을 다양한 방법으로 표현할 수 있다. \\
B & 집합의 개념을 이해하고, 집합을 다양한 방법으로 표현할 수 있다. \\
C & 집합의 개념을 알고, 집합을 표현할 수 있다. \\
D & 집합인 것과 아닌 것을 구분하고, 집합을 표현할 수 있다. \\
E & 집합인 것과 아닌 것을 구분할 수 있다. \\
\bottomrule
\end{tabularx}

\section{다음 차시}
\begin{calloutbox}
\textbf{02 집합 사이의 포함 관계 --- 1/1차시.} `01 집합'이 여기서 마무리된다. 다음 시간부터는 두 집합 사이의 관계(부분집합)를 다룬다.
\end{calloutbox}
\end{document}
